\begin{landscape}
\pdfbookmark[2]{First Calculation, [4]}{pallas-III-fragment4}
\begin{center}[4.]\end{center}
\begingroup
\makeatletter\tagsleft@true\let\veqno\@@leqno\makeatother
\setlength{\lineskip}{4pt}\setlength{\lineskiplimit}{4pt}
\setlength{\arraycolsep}{3pt}
\newcommand{\IIIfourDayComma}{\mathord{\rlap{\(^{\mathrm d}\)}{,}\kern.2em}}
\newcommand{\IIIfourVStroke}{\rlap{\kern.18em\raisebox{.6ex}{\rotatebox{25}{\rule{.25em}{.2pt}}}}}
\newcommand{\IIIfourPlusMark}{\mathord{+}\text{\rlap{\kern-.1em\raisebox{1.4ex}{\tiny .}}}}
\newcommand{\IIIfourLowComma}{\mathord{\text{\rlap{\kern.1em\raisebox{-.7ex}{.}},}}}
\newcommand{\IIIfourLowC}{\mathord{\underset{\smash{\raisebox{.5ex}{\tiny .}}}{c}}}
\noindent\hspace*{\parindent}[We now have to determine the values of the normal elements that, in combination with these perturbations, best represent the oppositions. For this purpose we shall follow in broad outline the method set out in Arts. 5--6 of the Disquisitio de elementis ellipticis Palladis (see also the example in Art. 8), where (Vol. VI, p. 4) the observations of the six oppositions are given as follows.
\begin{center}
\begingroup\setlength{\tabcolsep}{8pt}
\begin{tabular}{@{}c|r|r|r@{}}
Opposition & \(1803+\) & heliocentric longitude \(=\alpha\) & geocentric latitude \(=\beta\)\\
\hline
I & \(181\IIIfourDayComma019120\) & \(277^0\,39'\,24'',0\) & \(+46^0\,26'\,36'',0\)\\
II & \(608,207257\) & \(337\quad0\quad36,1\) & \(+15\quad1\quad49,8\)\\
III & \(1064,468796\) & \(67\quad20\quad42,9\) & \(-54\quad30\quad54,9\)\\
IV & \(1585,609502\) & \(223\quad37\quad27,7\) & \(\IIIfourPlusMark42\quad11\quad25,6\)\\
V\IIIfourVStroke & \(2034,887176\) & \(304\quad2\quad59,7\) & \(+37\quad43\quad53,7\)\\
VI & \(2457,673843\) & \(359\quad40\quad4,4\) & \(-7\quad22\quad10,1\)
\end{tabular}\endgroup
\end{center}
\noindent In that work the following notation was introduced.
\begin{center}
\SourceLeaf{486}
\begin{tabular}{@{}r@{\ }p{.90\linewidth}@{}}
\(\alpha={}\)&observed heliocentric longitude\\[3pt]
\(\beta={}\)&observed geocentric latitude\\[3pt]
\(\delta={}\)&heliocentric latitude, derived from \(\alpha\) by the formula \(\operatorname{tang}\delta=\operatorname{tang}i\sin(\alpha-\AscendingNode)\)\\[3pt]
\(\gamma={}\)&heliocentric latitude, derived from \(\beta\) by the formula \(\sin(\beta-\gamma)=\dfrac{R\sin\beta}{r}\).
\end{tabular}
\end{center}
\par
The following approximate values were assumed for the four elements \(i\), \(\AscendingNode\), \(\varphi\), \(\tilde\omega\), referred to the equinox 1803,0.
\[\begin{aligned}
i&=34^0\,37'\,47''&\qquad\AscendingNode&=172^0\,28'\,32''\\
\varphi&=14\quad12\quad51&\tilde\omega&=121\quad7\quad54.
\end{aligned}\]
\par
From the six observed heliocentric longitudes, reduced to the equinox just mentioned,
\[\begin{array}{@{}r@{\;}r|r|r|r|r|r@{}}
\alpha={}&277^{0}\,38'59'',17 & 336^{0}\,59'12'',67 & 67^{0}\,18'16'',89 & 223^{0}\,33'50'',20 & 303^{0}\,58'20'',58 & 359^{0}\,34'27'',29
\end{array}\]
and from the following values of \(i\) and \(\AscendingNode\) assumed for the six times with the perturbations taken into account,
\[\begin{array}{@{}r@{\;}r|r|r|r|r|r@{}}
\AscendingNode={}&172^{0}\,28'32'',43 & 28'49'',88 & 28'32'',77 & 28'19'',28 & 27'48'',22 & 27'45'',90
\\
i={}&34\ 37\ 46,96 & 37\ 35,85 & 37\ 12,42 & 37\ 25,19 & 37\ 30,77 & 36\ 39,79
\end{array}\]
we find, by the formulas
\[ \operatorname{tang}\delta = \operatorname{tang} i \sin(\alpha-\AscendingNode) \]
\[ \operatorname{tang}(v-\AscendingNode) = \frac{\operatorname{tang}(\alpha-\AscendingNode)}{\cos i} \]
the following heliocentric latitudes
\[\begin{array}{@{}r@{\;}r|r|r|r|r|r@{}}
\delta={}&+33^{0}\,41'6'',30 & +10^{0}\,27'5'',72 & -33^{0}\,40'33'',54 & +28^{0}\,14'54'',88 & +27^{0}\,20'33'',96 & -4^{0}\,53'0'',48
\end{array}\]
and the longitudes in the orbit
\[\begin{array}{@{}r@{\;}r|r|r|r|r|r@{}}
v={}&275^{0}\,3'20'',76 & 333^{0}\,51'47'',21 & 69^{0}\,53'48'',50 & 228^{0}\,52'49'',17 & 298^{0}\,31'39'',70 & 1^{0}\,4'56'',77.
\end{array}\]
\par
If we further set the longitude of perihelion and the angle \(\varphi\) according to the above perturbation values, as follows,
\[\begin{array}{@{}r@{\;}r|r|r|r|r|r@{}}
\tilde\omega={}&121^{0}\,7'53'',99 & 6'38'',57 & 3'17'',76 & 0'53'',65 & 120^{0}\,59'36'',29 & 45'59'',50
\\
\varphi={}&14\ 12\ 47,43 & 11\ 33,15 & 10\ 58,62 & 10\ 10,05 & 14\ \ 8\ 10,26 & 1\ 34,91
\end{array}\]
the true and mean anomalies and the mean longitudes follow.
\[\begin{array}{@{}r@{\;}r|r|r|r|r|r@{}}
v-\tilde\omega={}&153^{0}\,55'26'',77 & 212^{0}\,45'8'',64 & 308^{0}\,50'30'',74 & 104^{0}\,51'55'',52 & 177^{0}\,32'3'',41 & 240^{0}\,18'57'',27
\\
M={}&139\ 10\ 15,32 & 230\ 38\ 18,77 & 328\ 18\ 11,43 & 79\ 50\ 57,16 & 176\ \ 3\ 51,83 & 266\ 36\ \ 8,99
\\
L={}&260\ 18\ \ 9,31 & 351\ 44\ 57,34 & 89\ 21\ 29,19 & 200\ 51\ 50,81 & 297\ \ 3\ 28,12 & 27\ 22\ \ 8,49.
\end{array}\]
\par
Applying the perturbation quantities \(\delta\int n\,\mathrm{d}t+\delta\varepsilon\) to the latter values with their signs reversed gives the following mean longitudes free of perturbations.
\[\begin{array}{@{}r@{\;}r|r|r|r|r|r@{}}
L_0={}&260^{0}\,18'4'',58 & 351^{0}\,45'29'',31 & 89^{0}\,26'34'',76 & 201^{0}\,0'49'',51 & 297^{0}\,11'59'',74 & 27^{0}\,43'4'',39
\end{array}\]
Finally, we derive from these values the epoch of the mean longitude for \(1803.156^{\mathrm d}\) by subtracting the product of the mean motion, taken to be \(770'',7332\), and the intervening time. This gives
\[\begin{array}{@{}r@{\;}r|r|r|r|r|r@{}}
(\varepsilon_0)={}&254^{0}\,56'41'',51 & 56'38'',10 & 56'47'',58 & 56'42'',41 & 56'50'',65 & 56'47'',64.
\end{array}\]
\par
\SourceLeaf{487}These six values would be equal if our approximate values for the normal elements \(\tilde\omega\), \(\varphi\) and \(n\) were correct. From the differences among them we can now derive the corrections to the longitude of perihelion, the angle \(\varphi\), and the mean motion, together with the true value of the epoch of the mean longitude. If \(\Delta\tilde\omega\) and \(\Delta\varphi\) are the corrections to the quantities \(\tilde\omega\) and \(\varphi\), the corrected value of the mean longitude at each opposition is as follows (see Vol. VI, top of p. 7)
\[ L_0 + (m+1)\Delta\tilde\omega + n\Delta\varphi, \]
where the coefficients \(m\) and \(n\) have the same meaning as in the passage cited. If \(\Delta n\) is the correction to the mean motion, then the corrected epoch of the mean longitude, after subtracting the perturbations, is
\[ \varepsilon_0 = (\varepsilon_0) + (m+1)\Delta\tilde\omega + n\Delta\varphi - t\Delta n \tag*{A)}\]
where \(t\) denotes the time elapsed since \(1803.156^{\mathrm d}\) and \((\varepsilon_0)\) denotes the values calculated above from the observations.
\par
We can first eliminate the epoch from these six equations by taking their mean, which gives the corrected value of the epoch in the form
\[ \varepsilon_0 = \tfrac{1}{6}\Sigma(\varepsilon_0) + \tfrac{1}{6}\Sigma(m+1).\Delta\tilde\omega + \tfrac{1}{6}\Sigma n.\Delta\varphi - \tfrac{1}{6}\Sigma t.\Delta n \tag*{B)}\]
Subtracting each of the preceding equations A) from this mean B) gives six equations involving only the unknowns \(\Delta\tilde\omega\), \(\Delta\varphi\), \(\Delta n\). We solve these by the method of least squares. In the present case the numerical calculation shows that \(\Delta\tilde\omega\) and \(\Delta\varphi\) are extremely small, while
\[ \Delta n = +0'',0040 \]
follows, so that the mean motion becomes
\[ n = 770'',7372. \]
\par
Equation B) then gives the value of the epoch.
\[ \varepsilon_0 = 254^{0}\,56'40'',14. \]
\par
Correcting the six values of \((\varepsilon_0)\) found above by applying the corrections just obtained according to equations A) gives the following values.
\[\begin{array}{@{}r@{\;}r|r|r|r|r|r@{}}
\varepsilon_0={}&254^{0}\,56'41'',41 & 56'36'',35 & 56'44'',07 & 56'36'',88 & 56'43'',38 & 56'38'',74
\end{array}\]
Only the following very small deviations of the mean longitudes thus remain (see page 418 of this volume and Vol. VI, p. 320).
\[\begin{array}{@{}r@{\;}r|r|r|r|r|r@{}}
&+1'',27 & -3'',79 & +3'',93 & -3'',26 & +3'',24 & -1'',40.
\end{array}\]
\par
To improve the values of \(i\) and \(\AscendingNode\), we proceed as follows. The values of \(n\) just obtained give, for the several oppositions,
\[\begin{array}{@{}r@{\;}r|r|r|r|r|r@{}}
n={}&770'',7255 & 770'',3928 & 770'',1373 & 770'',5010 & 770'',1580 & 768''\IIIfourLowComma2969
\end{array}\]
from which we obtain
\[\begin{array}{@{}r@{\;}r|r|r|r|r|r@{}}
p={}&0{,}415\,0673 & 0{,}415\,2715 & 0{,}415\,4043 & 0{,}415\,3194 & 0{,}415\,5755 & 0{,}416\,6936
\end{array}\]
and
\[\begin{array}{@{}r@{\;}r|r|r|r|r|r@{}}
r={}&0{,}523\,2726 & 0{,}515\,5612 & 0{,}353\,3230 & 0{,}449\,2234 & 0{,}537\,0545 & 0{,}472\,2233
\end{array}\]
\SourceLeaf{488}
Then the formula \(\sin(\beta-\gamma) = \dfrac{R\sin\beta}{r}\) gives
\[\begin{array}{@{}r@{\;}r|r|r|r|r|r@{}}
\beta-\gamma={}&12^{0}\,45'32'',02 & 4^{0}\,34'41'',98 & 20^{0}\,50'28'',02 & 13^{0}\,56'26'',89 & 10^{0}\,23'37'',27 & 2^{0}\,29'4'',78
\\
\gamma={}&+33\ 41\ \ 3,98 & +10\ 27\ \ 7,82 & -33\ 40\ 26,88 & +28\ 14\ 58,71 & +27\ 20\ 16,43 & -4\ 53\ \ 5,32
\end{array}\]
and hence finally the difference
\[\begin{array}{@{}r@{\;}r|r|r|r|r|r@{}}
\gamma-\delta={}&-2'',32 & +2'',10 & +6'',66 & +3'',83 & -17'',53 & -4'',84.
\end{array}\]
\par
According to Art. 6 of the Disquisitio we then obtain for each opposition an equation of the form
\[ \IIIfourLowC\Delta\AscendingNode + f\Delta i = \gamma-\delta \]
where the quantities \(a\) and \(b\) have been omitted and \(c\) and \(f\) have nearly the same numerical values as in the passage cited. Our equations give the corrections if we set aside opposition V, since it was poorly observed.
\[ \Delta\AscendingNode = -2'',5 \quad\text{and}\quad \Delta i = -1'', \]
which gives
\[ \AscendingNode = 172^{0}\,28'29'',5 \quad\text{and}\quad i = 34^{0}\,37'46'' \]
After applying the resulting corrections to the above values of \(\gamma\), the following differences \(\gamma-\delta\) in the geocentric latitudes remain (see page 418 of this volume and Vol. VI, p. 320).
\[\begin{array}{@{}r@{\;}r|r|r|r|r|r@{}}
&-1'',00 & +4'',11 & +5'',99 & +3'',98 & -15'',86 & -3'',34.
\end{array}\]
\par
The resulting system of normal elements is as follows.
\[\begin{aligned}
n&=770'',7372\\
\tilde\omega&=121^{0}\,7'54''\\
\varphi&=14\ 12\ 51\\
\log a&=0{,}442\,0669\\
\varepsilon&=254^{0}\,56'40'',14\\
i&=34\ 37\ 46\\
\AscendingNode&=172\ 28\ 29,5.]
\end{aligned}\]
\par\medskip\begin{center}\rule{.45\linewidth}{.3pt}\end{center}
\par\endgroup
\end{landscape}
